\documentclass[UTF8, 12pt, fontset=fandol, oneside]{ctexart}
\usepackage{researchnotes}

% 保留分册独立编译时的章号前缀。
\renewcommand{\thesection}{12.\arabic{section}}

\title{第十二章\quad 傅里叶级数}
\author{}
\date{}

\begin{document}

\maketitle

在 18 世纪中叶, 法国数学家傅里叶 (Fourier) 在研究热传导问题时首先引进了三角级数, 后来人们将其称为傅里叶级数. 傅里叶级数理论 (人们后来称与之有关的理论为傅里叶分析) 不仅是数学研究中的一门重要的分支, 而且在物理学、工程等方面都有重要的应用. 在数学分析课程中, 我们主要介绍傅里叶级数的一些最基本性质。

我们可以从下面纯数学的观点来引进傅里叶级数。

在上一章中, 我们已经较系统地学习了幂级数. 现在换一种观点来看幂级数, 我们可以对幂级数的形式作如下的描述。

取定一组函数系, 即
\[
  1,\ x,\ x^2,\ \cdots,\ x^n,\ \cdots,
\]
则任何一个多项式 $P_n(x)$ 可以认为是从该函数系中取有限个元素的线性组合, 而一个幂级数 $\sum\limits_{n=0}^{+\infty}a_nx^n$ 则是该函数系的一个无穷的线性组合。

从这种观点来看, 我们可以取另一组函数系来考虑其线性组合. 要得到系统的理论, 必须要求这个函数系具有很好的性质以及其线性组合也具有很好的性质。

我们所熟悉的函数中, 除了多项式外, $\sin nx,\ \cos nx(n=1,2,\cdots)$ 则是遇到比较多的一类初等函数. 而三角函数系的无穷线性组合就是我们前面提到的傅里叶级数. 在本章中, 我们将考虑三角函数系的无穷线性组合——傅里叶级数, 并研究其基本性质。

\section{函数的傅里叶级数}

\subsection{基本三角函数系}

形如
\[
  \frac{a_0}{2}+\sum_{n=1}^{+\infty}(a_n\cos nx+b_n\sin nx)
\]
的级数称为傅里叶级数, 其中 $a_0,a_n,b_n(n=1,2,\cdots)$ 为实数. 我们可以将傅里叶级数看成函数系
\begin{equation}
  1,\ \cos x,\ \sin x,\ \cos 2x,\ \sin 2x,\ \cdots,\ \cos nx,\ \sin nx,\ \cdots
  \tag{12.1.1}
\end{equation}
的一个线性组合, 这个函数系我们以后称之为基本三角函数系. 为了以后叙述方便, 我们称基本三角函数系中有限个元素的线性组合为一个三角多项式. 特别地, 称
\[
  \frac{a_0}{2}+\sum_{k=1}^n(a_k\cos kx+b_k\sin kx)=T_n(x)
\]
为一个 $n$ 阶三角多项式。

为了讨论傅里叶级数, 首先我们来研究函数系 (12.1.1). 基本三角函数系 (12.1.1) 的第一个特性是它的周期性, 即基本函数系 (12.1.1) 中的函数具有公共的周期 $2\pi$。

我们对闭区间 $[a,b]$ 上的可积函数 $f(x),g(x)$, 定义它们之间的内积为
\[
  \int_a^b f(x)g(x)\mathrm dx.
\]
如果 $[a,b]$ 上的可积函数 $f(x),g(x)$ 满足 $\int_a^b f(x)g(x)\mathrm dx=0$, 则称 $f(x)$ 和 $g(x)$ 在 $[a,b]$ 上正交。

基本三角函数系 (12.1.1) 的第二个特性是它的正交性, 即基本函数系 (12.1.1) 中任意两个不同的函数在长度为 $2\pi$ 的任意区间上正交. 由于该函数系的周期性, 上述断言等价于下述等式成立:
\[
  (1)\quad \int_{-\pi}^{\pi}\sin mx\sin nx\mathrm dx=0
  \quad (m\ne n,\ \text{并且 }n,m\in\mathbb N);
\]
\[
  (2)\quad \int_{-\pi}^{\pi}\cos mx\cos nx\mathrm dx=0
  \quad (m\ne n\ \text{并且 }n,m\in\mathbb N\cup\{0\});
\]
\[
  (3)\quad \int_{-\pi}^{\pi}\cos mx\sin nx\mathrm dx=0
  \quad (m\in\mathbb N\cup\{0\},\ n\in\mathbb N).
\]
事实上, 对于 $\forall n,m\in\mathbb N$, 我们有
\[
\begin{aligned}
  \int_{-\pi}^{\pi}\sin mx\sin nx\mathrm dx
  &=\frac12\int_{-\pi}^{\pi}[\cos(m-n)x-\cos(m+n)x]\mathrm dx \\
  &=
  \begin{cases}
    \dfrac12\left[\dfrac{\sin(m-n)x}{m-n}-\dfrac{\sin(m+n)x}{m+n}\right]\bigg|_{-\pi}^{\pi}=0,
      & m\ne n,\\[1.2em]
    \dfrac12\left[x-\dfrac{\sin(m+n)x}{m+n}\right]\bigg|_{-\pi}^{\pi}=\pi,
      & m=n;
  \end{cases}
\end{aligned}
\]
\[
\begin{aligned}
  \int_{-\pi}^{\pi}\cos mx\cos nx\mathrm dx
  &=\frac12\int_{-\pi}^{\pi}[\cos(m-n)x+\cos(m+n)x]\mathrm dx \\
  &=
  \begin{cases}
    0, & m\ne n,\\
    \pi, & m=n;
  \end{cases}
\end{aligned}
\]
\[
  \int_{-\pi}^{\pi}\sin mx\cos nx\mathrm dx
  =\frac12\int_{-\pi}^{\pi}[\sin(m+n)x+\sin(m-n)x]\mathrm dx=0.
\]
另外, 我们有
\[
  \int_{-\pi}^{\pi}1^2\mathrm dx=2\pi
\]
和
\[
  \int_{-\pi}^{\pi}\sin nx\cdot1\mathrm dx
  =\int_{-\pi}^{\pi}\cos nx\cdot1\mathrm dx=0.
\]
因此 (1), (2) 和 (3) 成立。

\subsection{周期为 $2\pi$ 的函数的傅里叶级数}

与幂级数理论一样, 一个函数能否展成傅里叶级数是傅里叶级数理论中的一个重要问题。

现在设
\[
  f(x)=\frac{a_0}{2}+\sum_{n=1}^{+\infty}(a_n\cos nx+b_n\sin nx)
\]
在 $[-\pi,\pi]$ 上成立, 且右边的级数一致收敛到 $f(x)$. 显然, 此时 $f(x)$ 是一个 $[-\pi,\pi]$ 上的连续函数. 由函数项级数的一般性理论知 $f(x)$ 在 $[-\pi,\pi]$ 上的定积分可通过该级数逐项积分得到. 对于 $\forall k\geqslant0$, 利用基本三角函数系的正交性得
\[
\begin{aligned}
  \int_{-\pi}^{\pi}f(x)\cos kx\mathrm dx
  &=\int_{-\pi}^{\pi}\frac{a_0}{2}\cos kx\mathrm dx
    +\sum_{n=1}^{+\infty}a_n\int_{-\pi}^{\pi}\cos nx\cos kx\mathrm dx \\
  &\quad+\sum_{n=1}^{+\infty}b_n\int_{-\pi}^{\pi}\sin nx\cos kx\mathrm dx \\
  &=\pi a_k,
\end{aligned}
\]
因此有
\[
  a_k=\frac1\pi\int_{-\pi}^{\pi}f(x)\cos kx\mathrm dx,\quad k=0,1,2,\cdots.
\]
由于对于 $\forall k\geqslant1$, 有
\[
\begin{aligned}
  \int_{-\pi}^{\pi}f(x)\sin kx\mathrm dx
  &=\int_{-\pi}^{\pi}\frac{a_0}{2}\sin kx\mathrm dx
    +\sum_{n=1}^{+\infty}a_n\int_{-\pi}^{\pi}\cos nx\sin kx\mathrm dx \\
  &\quad+\sum_{n=1}^{+\infty}b_n\int_{-\pi}^{\pi}\sin nx\sin kx\mathrm dx \\
  &=\pi b_k,
\end{aligned}
\]
因此有
\[
  b_k=\frac1\pi\int_{-\pi}^{\pi}f(x)\sin kx\mathrm dx,\quad k=1,2,\cdots.
\]
以上的分析给我们提供了两点信息:

(1) 如果一个函数 $f(x)$ 能展成傅里叶级数并且其傅里叶级数能逐项积分, 则该 $f(x)$ 的傅里叶级数具有唯一性。

(2) 如果一个函数 $f(x)$ 在 $[-\pi,\pi]$ 上可积, 我们就可计算出
\begin{equation}
  a_n=\frac1\pi\int_{-\pi}^{\pi}f(x)\cos nx\mathrm dx,\quad n=0,1,2,\cdots
  \tag{12.1.2}
\end{equation}
及
\begin{equation}
  b_n=\frac1\pi\int_{-\pi}^{\pi}f(x)\sin nx\mathrm dx,\quad n=1,2,\cdots,
  \tag{12.1.3}
\end{equation}
并形式地得到一个与 $f(x)$ 有关系的傅里叶级数. 为此我们形式地将它记为
\begin{equation}
  f(x)\sim\frac{a_0}{2}+\sum_{n=1}^{+\infty}(a_n\cos nx+b_n\sin nx),
  \tag{12.1.4}
\end{equation}
其中 $a_n(n=0,1,2,\cdots), b_n(n=1,2,\cdots)$ 分别由 (12.1.2) 和 (12.1.3) 式给出. 今后我们将称 $a_n,b_n(n=1,2,\cdots)$ 为 $f(x)$ 的傅里叶系数, 并称
\[
  \frac{a_0}{2}+\sum_{n=1}^{+\infty}(a_n\cos nx+b_n\sin nx)
\]
为 $f(x)$ 的傅里叶级数。

在这里读者应注意到, 我们只是说 $f(x)$ 与该傅里叶级数有关 (在 (12.1.4) 式中, 我们用记号 $\sim$, 而不用等号 $=$), 这是因为我们并不知道该傅里叶级数是否能收敛于 $f(x)$. 我们很容易举出例子来说明该傅里叶级数很可能不点点收敛于 $f(x)$. 事实上, 改变 $f(x)$ 的有限个点的值, 并不改变 (12.1.2) 和 (12.1.3) 中的每个积分的值. 设 $f(x)$ 改变有限个点的值后得到的函数为 $g(x)$, 则 $f(x)$ 与 $g(x)$ 具有相同的傅里叶级数. 但是在这有限个点, 该傅里叶级数显然不可能同时收敛到 $f(x)$ 与 $g(x)$。

但对于连续函数 $f(x)$, 我们有以下结论。

\noindent\textbf{定理 12.1.1}\quad 设函数 $f(x)$ 在 $(-\infty,+\infty)$ 上连续, 以 $2\pi$ 为周期, 且 $f(x)$ 的傅里叶系数全为 $0$, 则 $f(x)$ 在 $(-\infty,+\infty)$ 上恒为 $0$。

\noindent\textbf{证明}\quad 倘若 $f(x)$ 在 $[-\pi,\pi]$ 上不恒为 $0$, 则在 $(-\pi,\pi)$ 内存在 $x_0$, 使得 $f(x_0)\ne0$. 我们不妨设 $f(x_0)>0$. 由 $f(x)$ 的连续性, 存在 $\delta>0$, 当 $x\in(x_0-\delta,x_0+\delta)$ 时, 有
\[
  f(x)>\frac{f(x_0)}{2}=M>0.
\]
由于 $f(x)$ 的傅里叶系数全为 $0$, 容易看出对于任何的三角多项式 $T(x)$, 有
\[
  \int_{-\pi}^{\pi}f(x)T(x)\mathrm dx=0.
\]
我们现在取定以下三角多项式
\[
\begin{aligned}
  T_0(x)&=1+\cos(x-x_0)-\cos\delta\\
  &=1+\cos x\cos x_0+\sin x\sin x_0-\cos\delta.
\end{aligned}
\]
显然存在 $r>1$, 对于 $x\in\left(x_0-\dfrac{\delta}{2},x_0+\dfrac{\delta}{2}\right)$, 有 $T_0(x)\geqslant r>1$. 注意到当 $x\in[x_0-\pi,x_0+\pi]\setminus(x_0-\delta,x_0+\delta)$ 时, 有 $|T_0(x)|\leqslant1$。

对于 $\forall n\in\mathbb N$, 由于 $T_0^n(x)$ 是一个三角多项式, 因此有
\begin{equation}
  \int_{x_0-\pi}^{x_0+\pi}f(x)T_0^n(x)\mathrm dx
  =\int_{-\pi}^{\pi}f(x)T_0^n(x)\mathrm dx=0.
  \tag{12.1.5}
\end{equation}
另一方面, 对于 $\forall n\in\mathbb N$, 我们有
\[
  \left|
    \int_{x_0-\pi}^{x_0-\delta}f(x)T_0^n(x)\mathrm dx
    +\int_{x_0+\delta}^{x_0+\pi}f(x)T_0^n(x)\mathrm dx
  \right|\leqslant 2\pi M\cdot1^n=2\pi M.
\]
而当 $n\to+\infty$ 时, 我们有
\[
  \int_{x_0-\delta}^{x_0+\delta}f(x)T_0^n(x)\mathrm dx
  \geqslant \int_{x_0-\frac{\delta}{2}}^{x_0+\frac{\delta}{2}}f(x)T_0^n(x)\mathrm dx
  \geqslant Mr^n\delta\to+\infty.
\]
因此有
\[
  \int_{x_0-\pi}^{x_0+\pi}f(x)T_0^n(x)\mathrm dx\to+\infty
  \quad (n\to+\infty).
\]
这与 (12.1.5) 矛盾. 证毕。

对于定理 12.1.1, 我们以后有更加简单的方法来证明之。

在举例来熟悉函数的傅里叶级数之前, 我们必须注意到以下事实: 对于在长度为 $2\pi$ 的半开半闭区间上定义的一个函数 $f(x)$, 我们总是可以以 $2\pi$ 为周期将它延拓到 $(-\infty,+\infty)$. 因此在下面的例子中的函数均是在一个长度为 $2\pi$ 的区间给出, 但它们可以看成是一个周期为 $2\pi$ 的函数的限制。

\noindent\textbf{例 12.1.1}\quad 设函数 $f(x)=x(-\pi\leqslant x<\pi)$, 求 $f(x)$ 的傅里叶级数。

\noindent\textbf{解}\quad $f(x)$ 周期延拓后的图像见图 12.1.1. 注意到 $\sin kx(k\in\mathbb N)$ 是奇函数, 而 $\cos kx(k\in\mathbb N\cup\{0\})$ 是偶函数, 因此, 当 $f(x)=x$ 为奇函数时, 有 $a_n=0(n=0,1,2,\cdots)$. 对于 $\forall n\in\mathbb N$, 我们有
\[
\begin{aligned}
  b_n
  &=\frac1\pi\int_{-\pi}^{\pi}x\sin nx\mathrm dx
  =\frac2\pi\int_0^\pi x\sin nx\mathrm dx\\
  &=-\frac2{n\pi}x\cos nx\bigg|_0^\pi+\frac2{n\pi}\int_0^\pi\cos nx\mathrm dx\\
  &=(-1)^{n-1}\frac2n.
\end{aligned}
\]
因此
\[
  f(x)\sim\sum_{n=1}^{+\infty}(-1)^{n-1}\frac2n\sin nx.
\]

在例 12.1.1 中, 严格说来 $f(x)$ 只在 $(-\pi,\pi)$ 中是奇函数. 但由于改变 $x=-\pi$ 一点处的函数值不会改变 $f(x)$ 的傅里叶系数, 因此我们不妨设 $f(x)$ 在 $[-\pi,\pi]$ 上是奇函数. 在下面的例子中, 我们说一个函数 $f(x)$ 是奇 (偶) 函数, 在大部分情形下仅是对 $f(x)$ 在一对称区间中除去若干个点后是奇 (偶) 函数而言的。

\noindent\textbf{思考题}\quad 设函数 $f(x)=x(x\in[0,2\pi))$, 是否 $f(x)$ 的傅里叶级数也是
\[
  \sum_{n=1}^{+\infty}(-1)^n\frac2n\sin nx?
\]

\noindent\textbf{例 12.1.2}\quad 设函数 $f(x)=\dfrac{\pi-x}{2}(x\in[0,2\pi))$, 求 $f(x)$ 的傅里叶级数。

\noindent\textbf{解}\quad 由图 12.1.2 可见, $f(x)$ 是奇函数, 因此 $a_n=0(n=0,1,2,\cdots)$. 对于 $\forall n\in\mathbb N$, 我们有
\[
\begin{aligned}
  b_n
  &=\frac2\pi\int_0^\pi\frac{\pi-x}{2}\sin nx\mathrm dx\\
  &=-\frac1{n\pi}(\pi-x)\cos nx\bigg|_0^\pi
    -\frac1{n\pi}\int_0^\pi\cos nx\mathrm dx\\
  &=\frac1n.
\end{aligned}
\]
因此我们有
\[
  f(x)\sim\sum_{n=1}^{+\infty}\frac{\sin nx}{n}.
\]

\noindent\textbf{注}\quad 从上例可以看出, 当考查 $f(x)$ 的奇偶性时, 可结合其延拓后的图像来考查, 不要仅仅从所给函数的解析式来下结论。

\noindent\textbf{例 12.1.3}\quad 设函数 $f(x)=x^2(x\in[0,2\pi))$, 求 $f(x)$ 的傅里叶级数。

\noindent\textbf{解}\quad $f(x)$ 周期延拓后的图像如图 12.1.3 所示。

从图 12.1.3 可以看出 $f(x)$ 不具有奇偶性. 由于 $f(x)$ 的傅里叶系数的积分中的被积函数均为以 $2\pi$ 为周期的函数, 因此我们只需在一个长度为 $2\pi$ 的区间来求之即可. 因此有
\[
  a_0=\frac1\pi\int_{-\pi}^{\pi}f(x)\mathrm dx
  =\frac1\pi\int_0^{2\pi}x^2\mathrm dx
  =\frac83\pi^2.
\]
\[
\begin{aligned}
  a_n
  &=\frac1\pi\int_0^{2\pi}x^2\cos nx\mathrm dx\\
  &=\frac1{n\pi}x^2\sin nx\bigg|_0^{2\pi}
    -\frac2{n\pi}\int_0^{2\pi}x\sin nx\mathrm dx\\
  &=\frac2{n^2\pi}x\cos nx\bigg|_0^{2\pi}
    -\frac2{n^2\pi}\int_0^{2\pi}\cos nx\mathrm dx\\
  &=\frac4{n^2},\quad n=1,2,\cdots,
\end{aligned}
\]
\[
\begin{aligned}
  b_n
  &=\frac1\pi\int_0^{2\pi}x^2\sin nx\mathrm dx\\
  &=-\frac1{n\pi}x^2\cos nx\bigg|_0^{2\pi}
    +\frac2{n\pi}\int_0^{2\pi}x\cos nx\mathrm dx\\
  &=-\frac{4\pi}{n}-\frac2{n^2\pi}\int_0^{2\pi}\sin nx\mathrm dx\\
  &=-\frac{4\pi}{n},\quad n=1,2,\cdots,
\end{aligned}
\]
因此
\[
  f(x)\sim\frac{4\pi^2}{3}
  +4\sum_{n=1}^{+\infty}\frac{\cos nx}{n^2}
  -4\pi\sum_{n=1}^{+\infty}\frac{\sin nx}{n}.
\]

\subsection{正弦级数与余弦级数}

若函数 $f(x)$ 在 $(a,b)$ 内有定义, 而且它是一个在 $[a,b]$ 上可积函数的限制, 我们将称 $f(x)$ 在 $(a,b)$ 内可积. 设 $f(x)$ 在 $(0,\pi)$ 内可积, 则我们可以将 $f(x)$ 延拓成 $(-\pi,\pi)$ 中的奇函数 $\widetilde f(x)$ (称之为奇延拓), 然后再将它周期地延拓到 $(-\infty,+\infty)$. 在这种情形下, $\widetilde f(x)$ 的傅里叶系数中
\[
  a_n=0,\quad n=0,1,2,\cdots,
\]
而
\[
  b_n=\frac1\pi\int_{-\pi}^{\pi}\widetilde f(x)\sin nx\mathrm dx
  =\frac2\pi\int_0^\pi f(x)\sin nx\mathrm dx,\quad n=1,2,\cdots,
\]
因此
\[
  f(x)\sim\sum_{n=1}^{+\infty}b_n\sin nx.
\]
此时我们称 $\sum\limits_{n=1}^{+\infty}b_n\sin nx$ 为 $f(x)$ 在 $(0,\pi)$ 上的正弦级数。

设 $f(x)$ 在 $(0,\pi)$ 内可积, 则我们也可以将 $f(x)$ 延拓成 $(-\pi,\pi)$ 中的偶函数 $\widetilde f(x)$ (称之为偶延拓), 然后再将它周期地延拓到 $(-\infty,+\infty)$. 在这种情形下, $\widetilde f(x)$ 的傅里叶系数中
\[
  b_n=0,\quad n=1,2,\cdots,
\]
而
\[
  a_n=\frac2\pi\int_0^\pi f(x)\cos nx\mathrm dx,\quad n=0,1,2,\cdots,
\]
因此
\[
  f(x)\sim\frac{a_0}{2}+\sum_{n=1}^{+\infty}a_n\cos nx.
\]
此时我们称 $\dfrac{a_0}{2}+\sum\limits_{n=1}^{+\infty}a_n\cos nx$ 为 $f(x)$ 在 $(0,\pi)$ 上的余弦级数. 由于在上述的奇延拓与偶延拓中, 延拓后的函数虽然不是唯一的, 但是仅仅在 $x=0,\pm\pi$ 处的值可能不同, 因此它们的傅里叶系数是相同的. 这说明了上面定义的正弦级数与余弦级数是由 $f(x)$ 在 $(0,\pi)$ 内的值唯一确定的。

\noindent\textbf{例 12.1.4}\quad 分别求函数 $f(x)=x^2(x\in(0,\pi))$ 的正弦级数与余弦级数。

\noindent\textbf{解}\quad 先求正弦级数. 由
\[
\begin{aligned}
  b_n
  &=\frac2\pi\int_0^\pi x^2\sin nx\mathrm dx\\
  &=-\frac2{n\pi}x^2\cos nx\bigg|_0^\pi
    +\frac4{n\pi}\int_0^\pi x\cos nx\mathrm dx\\
  &=(-1)^{n-1}\frac{2\pi}{n}
    +\frac4{n^2\pi}\left(x\sin nx\bigg|_0^\pi-\int_0^\pi\sin nx\mathrm dx\right)\\
  &=(-1)^{n-1}\frac{2\pi}{n}
    +\frac4{n^3\pi}[(-1)^n-1],\quad n=1,2,\cdots,
\end{aligned}
\]
我们求得 $f(x)$ 的正弦级数为
\[
  x^2\sim
  2\pi\sum_{n=1}^{+\infty}\frac{(-1)^n}{n}\sin nx
  -\frac8\pi\sum_{n=1}^{+\infty}\frac{\sin(2n-1)x}{(2n-1)^3}.
\]

再求余弦级数. 由
\[
  a_0=\frac2\pi\int_0^\pi x^2\mathrm dx=\frac23\pi^2,
\]
\[
\begin{aligned}
  a_n
  &=\frac2\pi\int_0^\pi x^2\cos nx\mathrm dx\\
  &=\frac2{n\pi}x^2\sin nx\bigg|_0^\pi
    -\frac4{n\pi}\int_0^\pi x\sin nx\mathrm dx\\
  &=\frac4{n^2\pi}\left(x\cos nx\bigg|_0^\pi-\int_0^\pi\cos nx\mathrm dx\right)\\
  &=(-1)^n\frac4{n^2},\quad n=1,2,\cdots,
\end{aligned}
\]
我们求得 $f(x)$ 的余弦级数为
\[
  x^2\sim\frac{\pi^2}{3}
  +4\sum_{n=1}^{+\infty}\frac{(-1)^n}{n^2}\cos nx.
\]

\subsection{周期为 $2T$ 的函数的傅里叶级数}

我们现在来讨论一下周期为 $2T>0$ 的函数的傅里叶级数问题。

设函数 $f(x)$ 以 $2T$ 为周期且在任何有限闭区间上可积, 作自变量变换 $x=\dfrac{T}{\pi}t$, 则函数 $F(t)=f\left(\dfrac{T}{\pi}t\right)$ 以 $2\pi$ 为周期且在任何有限闭区间上可积. 设
\[
  F(t)\sim\frac{a_0}{2}+\sum_{n=1}^{+\infty}(a_n\cos nt+b_n\sin nt),
\]
则将变量 $t$ 变回变量 $x$ 后, 我们便得到 $f(x)$ 的傅里叶级数:
\[
  f(x)\sim\frac{a_0}{2}
  +\sum_{n=1}^{+\infty}\left(a_n\cos\frac{n\pi}{T}x+b_n\sin\frac{n\pi}{T}x\right),
\]
其中
\[
  a_n=\frac1\pi\int_{-\pi}^{\pi}F(t)\cos nt\mathrm dt
  =\frac1T\int_{-T}^{T}f(x)\cos\frac{n\pi}{T}x\mathrm dx,\quad n=0,1,2,\cdots,
\]
\[
  b_n=\frac1\pi\int_{-\pi}^{\pi}F(t)\sin nt\mathrm dt
  =\frac1T\int_{-T}^{T}f(x)\sin\frac{n\pi}{T}x\mathrm dx,\quad n=1,2,\cdots.
\]

\noindent\textbf{例 12.1.5}\quad 求 $f(x)=x\cos x\left(-\dfrac{\pi}{2}\leqslant x<\dfrac{\pi}{2}\right)$ 的傅里叶级数。

\noindent\textbf{解}\quad $f(x)$ 是以 $\pi$ 为周期的函数, 且为奇函数, 因此 $a_n=0(n=0,1,2,\cdots)$. 记 $T=\dfrac{\pi}{2}$, 则
\[
\begin{aligned}
  b_n
  &=\frac2T\int_0^T x\cos x\sin 2nx\mathrm dx\\
  &=\frac1T\int_0^T x[\sin(2n-1)x+\sin(2n+1)x]\mathrm dx\\
  &=\frac1T\left[-\frac1{2n-1}x\cos(2n-1)x\bigg|_0^T
    +\frac1{2n-1}\int_0^T\cos(2n-1)x\mathrm dx\right]\\
  &\quad+\frac1T\left[-\frac1{2n+1}x\cos(2n+1)x\bigg|_0^T
    +\frac1{2n+1}\int_0^T\cos(2n+1)x\mathrm dx\right]\\
  &=\frac1T\left[\frac1{(2n-1)^2}\sin(2n-1)T
    +\frac1{(2n+1)^2}\sin(2n+1)T\right]\\
  &=\frac{(-1)^{n-1}}{T}\left[\frac1{(2n-1)^2}-\frac1{(2n+1)^2}\right]\\
  &=\frac{(-1)^{n-1}}{\pi}\frac{16n}{(4n^2-1)^2},\quad n=1,2,\cdots.
\end{aligned}
\]
因此
\[
  x\cos x\sim
  \sum_{n=1}^{+\infty}\frac{(-1)^{n-1}16n}{\pi(4n^2-1)^2}\sin 2nx.
\]

对于周期为 $2T$ 的函数 $f(x)$ 的傅里叶级数, 我们也可以用如下的办法来加以讨论. 取基本三角函数系
\[
  1,\ \cos\frac{\pi}{T}x,\ \sin\frac{\pi}{T}x,\ \cos\frac{2\pi}{T}x,\ \sin\frac{2\pi}{T}x,\ \cdots,\ 
  \cos\frac{n\pi}{T}x,\ \sin\frac{n\pi}{T}x,\ \cdots,
\]
容易看出它们是长度为 $2T$ 的区间上的正交函数系, 平行于周期为 $2\pi$ 的函数的傅里叶级数的讨论, 即可得到周期为 $2T$ 的 $f(x)$ 的傅里叶级数公式。

\section{傅里叶级数的敛散性}

从函数项级数 $\sum\limits_{n=1}^{+\infty}u_n(x)$ 的一般研究方法知, 若其前 $n$ 项部分和
\[
  S_n(x)=\sum_{k=1}^n u_k(x)\quad(n=1,2,\cdots)
\]
具有某种有规律的表达式, 则该级数可以容易通过对其部分和序列的研究来讨论其敛散性. 傅里叶级数正是具有这种特性, 由于它的部分和序列具有很有规律的表达式, 因此我们可以从它的研究中得到傅里叶级数敛散性的许多信息。

\subsection{狄利克雷积分}

从上节中我们已经知道, 若周期函数 $f(x)$ 在 $[-\pi,\pi]$ 上可积, 则可形式地得到其傅里叶级数
\[
  f(x)\sim\frac{a_0}{2}+\sum_{n=1}^{+\infty}(a_n\cos nx+b_n\sin nx),
\]
其中
\[
  a_n=\frac1\pi\int_{-\pi}^{\pi}f(x)\cos nx\mathrm dx,\quad n=0,1,2,\cdots
\]
\[
  b_n=\frac1\pi\int_{-\pi}^{\pi}f(x)\sin nx\mathrm dx,\quad n=1,2,\cdots.
\]
我们现在来考查 $f(x)$ 的傅里叶级数部分和序列. 对于 $\forall n\in\mathbb N$, 我们有
\[
\begin{aligned}
  S_n(x)
  &=\frac{a_0}{2}+\sum_{k=1}^n(a_k\cos kx+b_k\sin kx)\\
  &=\frac1{2\pi}\int_{-\pi}^{\pi}f(u)\mathrm du
    +\sum_{k=1}^n\frac1\pi\int_{-\pi}^{\pi}f(u)[\cos kx\cos ku+\sin kx\sin ku]\mathrm du\\
  &=\frac1\pi\int_{-\pi}^{\pi}f(u)\left[\frac12+\sum_{k=1}^n\cos k(u-x)\right]\mathrm du\\
  &=\frac1\pi\int_{-\pi}^{\pi}f(u)\frac{\sin\left(n+\frac12\right)(u-x)}{2\sin\frac{u-x}{2}}\mathrm du.
\end{aligned}
\]
为了更加便于分析, 我们作变换 $u=x+t$, 并利用 $f(x)$ 的周期性得
\[
\begin{aligned}
  S_n(x)
  &=\frac1\pi\int_{-\pi+x}^{\pi+x}
    f(u)\frac{\sin\left(n+\frac12\right)(u-x)}{2\sin\frac{u-x}{2}}\mathrm du\\
  &=\frac1\pi\int_{-\pi}^{\pi}
    f(x+t)\frac{\sin\left(n+\frac12\right)t}{2\sin\frac t2}\mathrm dt\\
  &=\frac1\pi\int_{-\pi}^{0}
    f(x+t)\frac{\sin\left(n+\frac12\right)t}{2\sin\frac t2}\mathrm dt
    +\frac1\pi\int_0^\pi
    f(x+t)\frac{\sin\left(n+\frac12\right)t}{2\sin\frac t2}\mathrm dt\\
  &=\frac1\pi\int_0^\pi
    (f(x+t)+f(x-t))\frac{\sin\left(n+\frac12\right)t}{2\sin\frac t2}\mathrm dt.
\end{aligned}
\]
上述 $f(x)$ 的傅里叶级数部分和的积分形式将给我们提供许多有用信息. 下面我们先来对它们作一些粗略的分析. 若取 $f(x)\equiv1$, 则其傅里叶级数的前 $n$ 项的部分和 $S_n(x)\equiv1$. 这说明了对于 $\forall n\in\mathbb N$, 有
\[
  1=\frac2\pi\int_0^\pi\frac{\sin\left(n+\frac12\right)t}{2\sin\frac t2}\mathrm dt
  =2\int_0^\pi D_n(t)\mathrm dt,
\]
其中
\[
  D_n(t)=\frac{\sin\left(n+\frac12\right)t}{2\pi\sin\frac t2}.
\]
今后我们将称
\[
  \int_0^\pi [f(x+t)+f(x-t)]D_n(t)\mathrm dt
\]
为 $f(x)$ 的狄利克雷积分, 并且称 $D_n(t)$ 为狄利克雷核。

现取定 $x_0\in[-\pi,\pi]$, 我们来研究 $f(x)$ 的傅里叶级数部分和序列 $\{S_n(x_0)\}$ 是否以 $S_0$ 为极限. 由于
\[
\begin{aligned}
  S_n(x_0)-S_0
  &=\int_0^\pi (f(x_0+t)+f(x_0-t)-2S_0)D_n(t)\mathrm dt\\
  &=\int_0^\pi\frac{f(x_0+t)+f(x_0-t)-2S_0}{2\pi\sin\frac t2}
    \sin\left(n+\frac12\right)t\mathrm dt\\
  &=\int_0^\pi G(t)\sin\left(n+\frac12\right)t\mathrm dt,
\end{aligned}
\]
其中
\begin{equation}
  G(t)=\frac{f(x_0+t)+f(x_0-t)-2S_0}{2\pi\sin\frac t2}.
  \tag{12.2.1}
\end{equation}
对上述积分我们粗略地观察一下可以看出, 当 $n$ 很大时, $\sin\left(n+\dfrac12\right)t$ 的值随着 $t$ 在 $[0,\pi]$ 上的变化, 不断地在 $x$ 轴的上下波动. 若 $G(t)$ 是个常数的话, 显然有
\[
  \int_0^\pi G(t)\sin\left(n+\frac12\right)t\mathrm dt\to0
  \quad (n\to\infty).
\]
若 $G(t)$ 不是常数, 但 $G(t)$ 可积, 当 $n\to\infty$ 时, 则由 $\sin\left(n+\dfrac12\right)t$ 在 $x$ 轴上下不断波动, 使得 $G(t)\sin\left(n+\dfrac12\right)t$ 的积分不断相抵, 因此上面的积分也应该趋于 $0$。

在以下行文中, 我们经常要用到以下假定: $f(x)$ 在 $[a,b]$ 上可积或有瑕点时绝对可积. 对此假定的确切意义是: $f(x)$ 在 $[a,b]$ 上或者是黎曼可积的, 或者具有有限个瑕点, 在 $[a,b]$ 中不含瑕点的任何闭子区间 $f(x)$ 均是黎曼可积的, 并且 $|f(x)|$ 在 $[a,b]$ 上的瑕积分是收敛的。

我们有下面的黎曼-勒贝格引理。

\noindent\textbf{引理 12.2.1 (黎曼-勒贝格引理)}\quad 设函数 $f(x)$ 在区间 $[a,b]$ 上可积或有瑕点时绝对可积, 则
\[
  \lim_{\lambda\to+\infty}\int_a^b f(x)\sin\lambda x\mathrm dx
  =\lim_{\lambda\to+\infty}\int_a^b f(x)\cos\lambda x\mathrm dx=0.
\]

\noindent\textbf{证明}\quad 我们分两种情形加以证明。

(1) $f(x)$ 在 $[a,b]$ 上可积。

由 $f(x)$ 在 $[a,b]$ 上可积, 从而它是有界的, 因此存在 $M>0$, 使得 $\forall x\in[a,b]$, 有 $|f(x)|\leqslant M$. 再由 $f(x)$ 的可积性, 对于 $\forall\varepsilon>0$, 存在 $[a,b]$ 的分割
\[
  a=x_0<x_1<\cdots<x_N=b,
\]
对于 $i=1,2,\cdots,N$, 记 $\Delta x_i=x_i-x_{i-1}$ 和
\[
  M_i=\sup_{x\in[x_{i-1},x_i]}\{f(x)\},\quad
  m_i=\inf_{x\in[x_{i-1},x_i]}\{f(x)\},
\]
则有
\[
  \sum_{i=1}^N(M_i-m_i)\Delta x_i<\frac{\varepsilon}{2}.
\]
现在我们取定一个满足上述条件的固定分割, 注意到对每个 $1\leqslant i\leqslant N$, 当 $x\in[x_{i-1},x_i]$ 时, 有
\[
  0\leqslant f(x)-m_i\leqslant M_i-m_i,
\]
因此有
\[
\begin{aligned}
  \left|\int_a^b f(x)\sin\lambda x\mathrm dx\right|
  &=\left|\sum_{i=1}^N\int_{x_{i-1}}^{x_i}f(x)\sin\lambda x\mathrm dx\right|\\
  &\leqslant\left|\sum_{i=1}^N\int_{x_{i-1}}^{x_i}(f(x)-m_i)\sin\lambda x\mathrm dx\right|
    +\left|\sum_{i=1}^N m_i\int_{x_{i-1}}^{x_i}\sin\lambda x\mathrm dx\right|\\
  &\leqslant\sum_{i=1}^N(M_i-m_i)\Delta x_i+\frac2\lambda\sum_{i=1}^N|m_i|\\
  &<\frac{\varepsilon}{2}+\frac{2NM}{\lambda}.
\end{aligned}
\]
由于这个分割是固定的, 因此存在 $X=\dfrac{4NM}{\varepsilon}$, 当 $\lambda>X$ 时, 有 $\dfrac{2NM}{\lambda}<\dfrac{\varepsilon}{2}$. 所以, 当 $\lambda>X$ 时, 有
\[
  \left|\int_a^b f(x)\sin\lambda x\mathrm dx\right|<\varepsilon.
\]
这就证明了当 $f(x)$ 在 $[a,b]$ 上可积时, 有
\[
  \lim_{\lambda\to+\infty}\int_a^b f(x)\sin\lambda x\mathrm dx=0.
\]

(2) $f(x)$ 在 $[a,b]$ 上有瑕点。

不妨设 $f(x)$ 只有一个瑕点 $c\in(a,b)$. 由于 $f(x)$ 绝对可积, 因此存在 $\delta>0$, 使得 $a<c-\delta<c+\delta<b$ 和 $\int_{c-\delta}^{c+\delta}|f(x)|\mathrm dx<\dfrac{\varepsilon}{3}$, 从而对任意的 $\lambda$, 有
\[
  \left|\int_{c-\delta}^{c+\delta}f(x)\sin\lambda x\mathrm dx\right|
  \leqslant\int_{c-\delta}^{c+\delta}|f(x)|\mathrm dx<\frac{\varepsilon}{3}.
\]
再由于 $f(x)$ 在 $[a,c-\delta]$ 及 $[c+\delta,b]$ 上可积, 因此存在 $X>0$, 当 $\lambda>X$ 时, 有
\[
  \left|\int_a^{c-\delta}f(x)\sin\lambda x\mathrm dx\right|<\frac{\varepsilon}{3}
\]
及
\[
  \left|\int_{c+\delta}^b f(x)\sin\lambda x\mathrm dx\right|<\frac{\varepsilon}{3},
\]
从而当 $\lambda>X$ 时, 有
\[
\begin{aligned}
  \left|\int_a^b f(x)\sin\lambda x\mathrm dx\right|
  &\leqslant\left|\int_a^{c-\delta}f(x)\sin\lambda x\mathrm dx\right|
    +\left|\int_{c-\delta}^{c+\delta}f(x)\sin\lambda x\mathrm dx\right|\\
  &\quad+\left|\int_{c+\delta}^b f(x)\sin\lambda x\mathrm dx\right|\\
  &<\frac{\varepsilon}{3}+\frac{\varepsilon}{3}+\frac{\varepsilon}{3}=\varepsilon.
\end{aligned}
\]
因此由 (1) 和 (2) 推出在定理的假设下有
\[
  \lim_{\lambda\to+\infty}\int_a^b f(x)\sin\lambda x\mathrm dx=0.
\]
同理可证
\[
  \lim_{\lambda\to+\infty}\int_a^b f(x)\cos\lambda x\mathrm dx=0.
\]
证毕。

现在我们回过头来考查以下等式:
\[
  S_n(x_0)-S_0
  =\int_0^\pi\frac{f(x_0+t)+f(x_0-t)-2S_0}{2\pi\sin\frac t2}
    \sin\left(n+\frac12\right)t\mathrm dt.
\]
设 $f(x)$ 在 $[-\pi,\pi]$ 上可积或有瑕点时绝对可积, 我们来看该式当 $n\to\infty$ 时的变化情况。

在 $t=0$ 的某个邻域外, 若 $G(t)$ ($G(t)$ 由 (12.2.1) 定义) 没有瑕点显然是可积的, 并且有瑕点时也绝对可积, 因此由黎曼-勒贝格引理 (引理 12.2.1) 有
\[
  \int_\delta^\pi G(t)\sin\left(n+\frac12\right)t\mathrm dt\to0
  \quad (n\to\infty).
\]
所以当 $n$ 趋于无穷时, $S_n(x_0)-S_0$ 是否趋于 $0$ 仅与 $f(x_0\pm t)$ 在 $t=0$ 附近的值有关, 换句话说仅与 $f(x)$ 在 $x_0$ 附近的值有关. 因此我们有

\noindent\textbf{定理 12.2.2 (黎曼局部化定理)}\quad 设周期为 $2\pi$ 的函数 $f(x)$ 在 $[-\pi,\pi]$ 上可积或有瑕点时绝对可积, 则 $f(x)$ 的傅里叶级数在 $x_0\in[-\pi,\pi]$ 处的敛散性只与 $f(x)$ 在 $(x_0-\delta,x_0+\delta)$ 的值有关, 其中 $\delta>0$ 是一任意给定的正常数。

值得指出的是, 一个函数 $f(x)$ 的傅里叶系数是由 $f(x)$ 与基本三角函数系中每个函数的乘积在 $[-\pi,\pi]$ 上的积分得到. 换句话说, 它们依赖于 $[-\pi,\pi]$ 上 $f(x)$ 的取值. 而以上定理却告诉我们: 傅里叶级数在 $x_0$ 处的敛散性只依赖于 $f(x)$ 在 $x_0$ 附近的局部性质. 这个结果有点出乎人们的意料。

现在我们继续来研究 $f(x)$ 的傅里叶级数在 $x_0$ 处的敛散性. 为了更加容易估计 $S_n(x_0)-S_0$, 我们注意到在 $f(x)$ 的狄利克雷积分中被积函数的 $2\sin\dfrac t2$ 可以用 $t$ 代替. 换句话说, 若 $f(x)$ 在 $[0,\delta]$ 上可积或具有瑕点时绝对可积, 则当 $n\to\infty$ 时,
\[
  \int_0^\delta\frac{\sin\left(n+\frac12\right)t}{2\sin\frac t2}f(t)\mathrm dt
  \quad\text{与}\quad
  \int_0^\delta\frac{\sin\left(n+\frac12\right)t}{t}f(t)\mathrm dt
\]
具有相同的敛散性, 并且它们收敛时具有相同的极限。

事实上, 由
\[
  \lim_{t\to0}\left(\frac1{2\sin\frac t2}-\frac1t\right)
  =\lim_{t\to0}\frac{t-2\sin\frac t2}{2t\sin\frac t2}
  =\lim_{t\to0}\frac{t-2\left(\frac t2+o(t^2)\right)}{2t\sin\frac t2}=0,
\]
定义
\[
  \varphi(t)=
  \begin{cases}
    \dfrac1{2\sin\frac t2}-\dfrac1t, & t\ne0,\\
    0, & t=0,
  \end{cases}
\]
则 $\varphi(t)$ 在 $[0,\pi]$ 上连续, 因此有界, 从而 $f(t)\varphi(t)$ 在 $[0,\delta]$ 上可积或有瑕点时绝对可积. 由黎曼-勒贝格引理有
\[
  \lim_{n\to\infty}\int_0^\delta f(t)\varphi(t)\sin\left(n+\frac12\right)t\mathrm dt=0.
\]
由此即得
\[
  \lim_{n\to\infty}\int_0^\delta f(t)\frac{\sin\left(n+\frac12\right)t}{2\sin\frac t2}\mathrm dt
  =
  \lim_{n\to\infty}\int_0^\delta\frac{f(t)\sin\left(n+\frac12\right)t}{t}\mathrm dt.
\]

\noindent\textbf{注}\quad 上述等式的确切意义是, 等式两边的极限具有相同的敛散性, 并且它们收敛时具有相同的极限. 另外, 由黎曼-勒贝格引理及上面的分析, 若 $f(x)$ 在 $[0,\pi]$ 上可积或有瑕点时绝对可积, 则有
\[
  \lim_{n\to\infty}\int_0^\pi f(t)\frac{\sin\left(n+\frac12\right)t}{2\sin\frac t2}\mathrm dt
  =
  \lim_{n\to\infty}\int_0^\pi f(t)\frac{\sin\left(n+\frac12\right)t}{t}\mathrm dt.
\]
利用上述等式, 我们可以计算出积分 $\int_0^{+\infty}\dfrac{\sin x}{x}\mathrm dx$。

事实上, 我们已经知道该积分是收敛的, 取 $f(t)\equiv1$, 则有
\[
\begin{aligned}
  1
  &\equiv\lim_{n\to\infty}\frac2\pi\int_0^\pi
    \frac{\sin\left(n+\frac12\right)t}{2\sin\frac t2}\mathrm dt\\
  &=\lim_{n\to\infty}\frac2\pi\int_0^\pi
    \frac{\sin\left(n+\frac12\right)t}{t}\mathrm dt\\
  &=\lim_{n\to\infty}\frac2\pi\int_0^{(n+\frac12)\pi}\frac{\sin x}{x}\mathrm dx
  =\frac2\pi\int_0^{+\infty}\frac{\sin x}{x}\mathrm dx.
\end{aligned}
\]
因此
\[
  \int_0^{+\infty}\frac{\sin x}{x}\mathrm dx=\frac{\pi}{2}.
\]

\subsection{傅里叶级数的收敛判别法}

对于给定的 $x_0\in[-\pi,\pi]$, 我们继续进行周期为 $2\pi$ 的函数 $f(x)$ 的傅里叶级数在 $x_0$ 处敛散性的讨论. 根据上一小节推导, 我们知道 $f(x)$ 的傅里叶级数在 $x_0$ 处收敛到 $S_0$ 的充分必要条件是: 对充分小的正数 $\delta>0$, 有
\begin{equation}
  \lim_{n\to\infty}\int_0^\delta
  \frac{f(x_0+t)+f(x_0-t)-2S_0}{t}
  \sin\left(n+\frac12\right)t\mathrm dt=0.
  \tag{12.2.2}
\end{equation}
该极限的收敛性还可转化为如何选择 $S_1$ 及 $S_2$, 使得
\begin{equation}
  \lim_{n\to\infty}\int_0^\delta
  \frac{f(x_0+t)-S_1}{t}\sin\left(n+\frac12\right)t\mathrm dt=0
  \tag{12.2.3}
\end{equation}
及
\begin{equation}
  \lim_{n\to\infty}\int_0^\delta
  \frac{f(x_0-t)-S_2}{t}\sin\left(n+\frac12\right)t\mathrm dt=0.
  \tag{12.2.4}
\end{equation}
显然当 (12.2.3) 及 (12.2.4) 成立时, 令 $S_0=\dfrac{S_1+S_2}{2}$, 则 (12.2.2) 必成立。

我们首先来看一个简单的情形: 设 $f(x)$ 在 $x_0$ 处可导, 则
\[
  \lim_{t\to0}\frac{f(x_0+t)-f(x_0)}{t}=f'(x_0).
\]
特别地, 我们推知\allowbreak $\dfrac{f(x_0+t)-f(x_0)}{t}$ 在 $x_0$ 的邻域内有界. 这说明 $t=0$ 不是 $\dfrac{f(x_0+t)-f(x_0)}{t}$ 的瑕点, 因此取 $S_1=f(x_0)$, 则 (12.2.3) 成立. 同理, 当 $f(x)$ 在 $x_0$ 处可导时, 取 $S_2=f(x_0)$, 则 (12.2.4) 成立. 因此, 当 $f(x)$ 在 $x_0$ 处可导时, 取 $S_0=f(x_0)$, 则 (12.2.2) 成立。

基于上述的结果, 我们进一步讨论傅里叶级数在一个区间上的敛散性. 为此我们引入一个函数分段可微的概念。

设函数 $f(x)$ 在区间 $[a,b]$ 上有定义. 若存在 $[a,b]$ 的分割
\[
  \Delta: a=x_0<x_1<\cdots<x_n=b,
\]
使得 $f(x)$ 仅以 $x_i$ 为第一类间断点, 并且在 $x_i(i=1,2,\cdots,n)$ 处, 存在推广的单侧导数, 即
\[
  \lim_{h\to0+0}\frac{f(x_i+h)-f(x_i+0)}{h}\triangleq f'_+(x_0)
\]
与
\[
  \lim_{h\to0-0}\frac{f(x_i+h)-f(x_i-0)}{h}\triangleq f'_-(x_0)
\]
存在 (在 $[a,b]$ 的端点处存在推广的左右导数); 而当 $x\in(x_{i-1},x_i)(i=1,2,\cdots,n)$ 时, $f'(x)$ 存在. 若 $f(x)$ 满足以上条件, 则称 $f(x)$ 在 $[a,b]$ 上是分段可微的。

对于分段可微函数, 若在 (12.2.3) 式左边取 $S_1=f(x_0+0)$, 在 (12.2.4) 式左边取 $S_2=f(x_0-0)$, 则 (12.2.3), (12.2.4) 式同时成立. 因此我们有

\noindent\textbf{定理 12.2.3}\quad 设 $f(x)$ 是周期为 $2\pi$ 的函数, 且在 $[-\pi,\pi]$ 内分段可微, 则 $f(x)$ 的傅里叶级数处处收敛到
\[
  \frac{f(x-0)+f(x+0)}{2},
\]
即
\[
  \frac{a_0}{2}+\sum_{n=1}^{+\infty}(a_n\cos nx+b_n\sin nx)
  =
  \frac{f(x-0)+f(x+0)}{2},
  \quad x\in[-\pi,\pi].
\]

下面我们来看一些更为精细的结果. 对于 $x_0\in[-\pi,\pi]$, 我们令
\begin{equation}
  \varphi(t)=f(x_0+t)+f(x_0-t)-2S_0,
  \tag{12.2.5}
\end{equation}
则 $\varphi(t)$ 在 $t=0$(即 $f(x)$ 在 $x_0$) 附近的性质对傅里叶级数的收敛起着至关重要的作用。

由黎曼-勒贝格引理, 若要 (12.2.2) 成立, 只要 $\dfrac{\varphi(t)}{t}$ 在 $t=0$ 的邻域绝对可积即可. 因此我们有

\noindent\textbf{定理 12.2.4 (狄尼定理)}\quad 设 $f(x)$ 是周期为 $2\pi$ 的函数, 在 $[-\pi,\pi]$ 上可积或有瑕点时绝对可积, 并且对于 $x_0\in[-\pi,\pi]$, 存在 $\delta>0$, 使得
\[
  \int_0^\delta\frac{|\varphi(t)|}{t}\mathrm dt<+\infty,
\]
则 $f(x)$ 的傅里叶级数在 $x_0$ 处收敛到 $S_0$(其中 $\varphi(t)$ 见 (12.2.5))。

\noindent\textbf{证明}\quad 由定理所给条件知 $\dfrac{\varphi(t)}{(t)}$ 在 $[0,\delta]$ 上绝对可积. 又由 $f(x)$ 的条件知 $\dfrac{\varphi(t)}{(t)}$ 在 $[\delta,\pi]$ 上可积或有瑕点时绝对可积, 因此 $\dfrac{\varphi(t)}{t}$ 在 $[0,\pi]$ 上绝对可积. 由黎曼-勒贝格引理推出
\[
  \lim_{n\to\infty}[S_n(x_0)-S_0]
  =\lim_{n\to\infty}\int_0^\pi
  \frac{\varphi(t)}{t}\sin\left(n+\frac12\right)t\mathrm dt=0.
\]
证毕。

从定理 12.2.4 立即可推出下面的李普西茨 (Lipschitz) 定理。

\noindent\textbf{定理 12.2.5 (李普西茨定理)}\quad 设 $f(x)$ 是周期为 $2\pi$ 的函数, 在 $[-\pi,\pi]$ 上瑕可积或有瑕点时绝对可积, 再设 $f(x)$ 在 $x_0\in[-\pi,\pi]$ 处满足 $\alpha$-赫尔德 ($\alpha>0$) 连续性, 即存在 $L>0,\delta>0$, 使得对于 $x\in U(x_0,\delta)$, 有
\[
  |f(x_0+t)-f(x_0)|\leqslant L|t|^\alpha,
\]
则 $f(x)$ 的傅里叶级数在 $x_0$ 处收敛到 $f(x_0)$。

\noindent\textbf{证明}\quad 由 $f(x)$ 在 $x_0$ 处满足 $\alpha$-赫尔德 ($\alpha>0$) 连续性, 存在 $\delta>0$, 使得当 $t\in(-\delta,\delta)$ 时, 我们有
\[
  \frac{|f(x_0+t)-f(x_0)|}{t}\leqslant\frac{L}{t^{1-\alpha}}.
\]
由于 $\alpha>0$, 我们推出
\[
  \int_0^\delta\frac{|\varphi(t)|}{t}\mathrm dt<+\infty,
\]
其中 $\varphi(t)$ 见 (12.2.5). 由狄尼定理知 $f(x)$ 的傅里叶级数在 $x_0$ 处收敛到 $f(x_0)$. 证毕。

最后我们来证明一个在理论上及应用上都十分有用的结果, 即单调函数的傅里叶级数必定收敛。

\noindent\textbf{定理 12.2.6 (狄利克雷定理)}\quad 设 $f(x)$ 是周期为 $2\pi$ 的函数, 在 $[-\pi,\pi]$ 上可积或有瑕点时绝对可积, 再设 $x_0\in[-\pi,\pi]$ 不是 $f(x)$ 的瑕点, 并且存在 $\delta_0>0$, 使得 $f(x)$ 在 $(x_0-\delta_0,x_0)$ 及 $(x_0,x_0+\delta_0)$ 内分别单调, 则 $f(x)$ 的傅里叶级数在 $x_0$ 处收敛到
\[
  \frac{f(x_0-0)+f(x_0+0)}{2}.
\]

\noindent\textbf{证明}\quad 先设 $f(x)$ 在 $(x_0,x_0+\delta_0)$ 上单调. 我们证明
\begin{equation}
  \lim_{\lambda\to+\infty}\int_0^{\delta_0}
  \frac{f(x_0+t)-f(x_0+0)}{t}\sin\lambda t\mathrm dt=0.
  \tag{12.2.6}
\end{equation}
不妨设 $f(x)$ 是单调上升的. 由于
\[
  \int_0^{+\infty}\frac{\sin t}{t}\mathrm dt=\frac{\pi}{2},
\]
故连续函数 $G(x)=\displaystyle\int_0^x\dfrac{\sin t}{t}\mathrm dt$ 在 $[0,+\infty)$ 上有界. 由此存在常数 $M>1$, 使得对任意 $0\leqslant t_1<t_2<+\infty$, 有
\[
  \left|\int_{t_1}^{t_2}\frac{\sin t}{t}\mathrm dt\right|\leqslant M.
\]
由 $f(x)$ 的单调性, 对于 $\forall\varepsilon>0$, $\exists 0<\delta<\delta_0$, 使得当 $0<t<\delta$ 时, 有
\[
  0\leqslant f(x_0+t)-f(x_0+0)<\frac{\varepsilon}{2M}.
\]
为了估计积分, 我们记
\[
\begin{aligned}
  \int_0^{\delta_0}\frac{f(x_0+t)-f(x_0+0)}{t}\sin\lambda t\mathrm dt
  &=
  \int_0^\delta\frac{f(x_0+t)-f(x_0+0)}{t}\sin\lambda t\mathrm dt\\
  &\quad+
  \int_\delta^{\delta_0}\frac{f(x_0+t)-f(x_0+0)}{t}\sin\lambda t\mathrm dt\\
  &\triangleq I+J.
\end{aligned}
\]

对于 $J$, 注意到 $\dfrac{f(x_0+t)-f(x_0+0)}{t}$ 在 $[\delta,\delta_0]$ 上可积或有瑕点时绝对可积, 因此由黎曼-勒贝格引理知, $\exists\lambda_0>0$, 当 $\lambda>\lambda_0$ 时, 有
\[
  |J|<\frac{\varepsilon}{2}.
\]

现在来估计积分 $I$. 由定积分第二中值定理, 存在 $\xi\in(0,\delta)$, 使得
\[
\begin{aligned}
  |I|
  &=
  \left|\int_0^\delta\frac{f(x_0+t)-f(x_0+0)}{t}\sin\lambda t\mathrm dt\right|\\
  &=
  |f(x_0+\delta)-f(x_0+0)|
  \left|\int_\xi^\delta\frac1t\sin\lambda t\mathrm dt\right|\\
  &=
  |f(x_0+\delta)-f(x_0+0)|
  \left|\int_{\lambda\xi}^{\lambda\delta}\frac{\sin t}{t}\mathrm dt\right|\\
  &<
  \frac{\varepsilon}{2M}M=\frac{\varepsilon}{2}.
\end{aligned}
\]
因此
\[
  \left|\int_0^{\delta_0}\frac{f(x_0+t)-f(x_0+0)}{t}\sin\lambda t\mathrm dt\right|
  \leqslant |I|+|J|<\frac{\varepsilon}{2}+\frac{\varepsilon}{2}=\varepsilon.
\]
(12.2.6) 得证. 同理可证当 $f(x)$ 在 $[0,\delta_0]$ 上单调下降时, (12.2.6) 也成立。

由于 $f(x_0+t)$ 和 $f(x_0-t)$ 在 $[0,\delta_0]$ 上分别单调, 利用上面已证结果我们有
\[
  \lim_{n\to\infty}\int_0^{\delta_0}
  \left[
    f(x_0+t)+f(x_0-t)-2\frac{f(x_0+0)+f(x_0-0)}{2}
  \right]
  \frac{\sin\left(n+\frac12\right)t}{t}\mathrm dt=0.
\]
证毕。

\noindent\textbf{注}\quad 在定理的假设下, 由黎曼-勒贝格引理我们有
\[
  \lim_{n\to\infty}\int_{\delta_0}^\pi
  \left[
    f(x_0+t)+f(x_0-t)-2\frac{f(x_0+0)+f(x_0-0)}{2}
  \right]
  \frac{\sin\left(n+\frac12\right)t}{t}\mathrm dt=0,
\]
因此
\[
  \lim_{n\to\infty}\int_0^\pi
  \frac{f(x_0+t)+f(x_0-t)}{2}
  \frac{\sin\left(n+\frac12\right)t}{t}\mathrm dt
  =
  \frac{f(x_0+0)+f(x_0-0)}{2}.
\]

另外, 从定理 12.2.6 的证明中我们可以看出, 当 $f(x)$ 在 $x_0$ 的 $\delta_0$ 邻域内单调时, 下述等式成立:
\[
  \lim_{\lambda\to+\infty}\int_0^\pi f(x_0\pm t)\frac{\sin\lambda t}{t}\mathrm dt
  =\frac{\pi}{2}f(x_0\pm0).
\]

最后我们来小结一下本节的内容. 在实际应用中, 下述的结论是经常用到的。

设 $f(x)$ 是周期为 $2\pi$ 的函数, 若 $f(x)$ 满足以下条件之一:

(1) $f(x)$ 在 $[-\pi,\pi]$ 分段单调;

(2) $f(x)$ 在 $[-\pi,\pi]$ 分段可微,

则对于 $\forall x\in[-\pi,\pi]$, $f(x)$ 的傅里叶级数收敛到
\[
  \frac{f(x-0)+f(x+0)}{2}.
\]

因此, 对 \S 12.1 中的例子, 我们有

\noindent\textbf{例 12.2.1}\quad 函数 $f(x)=x(-\pi\leqslant x<\pi)$ 的傅里叶级数为
\[
  \sum_{n=1}^{+\infty}(-1)^{n-1}\frac2n\sin nx
  =
  \begin{cases}
    x, & x\in(-\pi,\pi),\\
    0, & x=-\pi,\pi.
  \end{cases}
\]

\noindent\textbf{例 12.2.2}\quad 函数 $f(x)=\dfrac{\pi-x}{2}(x\in[0,2\pi))$ 的傅里叶级数为
\[
  \sum_{n=1}^{+\infty}\frac{\sin nx}{n}
  =
  \begin{cases}
    \dfrac{\pi-x}{2}, & x\in(0,2\pi),\\
    0, & x=0,2\pi.
  \end{cases}
\]

\noindent\textbf{例 12.2.3}\quad 函数 $f(x)=x^2(x\in[0,2\pi))$ 的傅里叶级数为
\[
  \frac{4\pi^2}{3}+4\sum_{n=1}^{+\infty}\frac{\cos nx}{n^2}
  -4\pi\sum_{n=1}^{+\infty}\frac{\sin nx}{n}
  =
  \begin{cases}
    x^2, & x\in(0,2\pi),\\
    2\pi^2, & x=0,2\pi.
  \end{cases}
\]

\noindent\textbf{例 12.2.4}\quad 函数 $f(x)=x^2(x\in[0,\pi])$ 的正弦级数与余弦级数分别为
\[
  2\pi\sum_{n=1}^{+\infty}\frac{(-1)^n}{n}\sin nx
  -\frac8\pi\sum_{n=1}^{+\infty}\frac{\sin(2n-1)x}{(2n-1)^3}
  =
  \begin{cases}
    x^2, & x\in[0,\pi),\\
    0, & x=\pi
  \end{cases}
\]
和
\[
  \frac{\pi^2}{3}+4\sum_{n=1}^{+\infty}\frac{(-1)^n}{n^2}\cos nx=x^2,\quad x\in[0,\pi].
\]

\noindent\textbf{例 12.2.5}\quad 函数 $f(x)=x\cos x\left(-\dfrac{\pi}{2}\leqslant x<\dfrac{\pi}{2}\right)$ 的傅里叶级数为
\[
  \sum_{n=1}^{+\infty}\frac{(-1)^{n-1}}{\pi}\frac{16n}{(4n^2-1)^2}\sin2nx
  =x\cos x,\quad x\in\left[-\frac{\pi}{2},\frac{\pi}{2}\right].
\]

\noindent\textbf{注}\quad 在例 12.2.5 中函数的周期为 $\pi$, 但相应的傅里叶级数收敛性的结论仍真, 请读者自己证明该结论。

\noindent\textbf{例 12.2.6}\quad 设 $0<\alpha<1$, 试求函数 $f(x)=\cos\alpha x(x\in(-\pi,\pi))$ 的傅里叶级数, 并求出该傅里叶级数的和函数。

\noindent\textbf{解}\quad 显然, $f(x)$ 为偶函数, 因此
\[
  b_n=0,\quad n=1,2,\cdots;
\]
\[
  a_0=\frac2\pi\int_0^\pi\cos\alpha x\mathrm dx=\frac{2}{\pi\alpha}\sin\alpha\pi,
\]
\[
\begin{aligned}
  a_n
  &=\frac2\pi\int_0^\pi\cos\alpha x\cos nx\mathrm dx\\
  &=\frac1\pi\int_0^\pi[\cos(\alpha-n)x+\cos(\alpha+n)x]\mathrm dx\\
  &=\frac1\pi\left[
    \frac{\sin(\alpha-n)x}{\alpha-n}
    +\frac{\sin(\alpha+n)x}{\alpha+n}
  \right]_0^\pi\\
  &=(-1)^n\frac{2\alpha\sin\alpha\pi}{\pi(\alpha^2-n^2)},\quad n=1,2,\cdots.
\end{aligned}
\]
注意到 $\cos\alpha x$ 分段可微且连续, 因此当 $x\in[-\pi,\pi]$ 时, 有
\[
  \frac{\sin\alpha\pi}{\pi}\left(
    \frac1\alpha+\sum_{n=1}^{+\infty}\frac{(-1)^n2\alpha}{\alpha^2-n^2}\cos nx
  \right)=\cos\alpha x.
\]

利用以上展开式, 我们可以得到一些有趣的推论 (见本章习题). 最后我们简单地介绍一下关于傅里叶级数敛散性的研究历史. 傅里叶在研究热传导问题时引进了傅里叶级数后, 当时他曾认为基本上所有函数的傅里叶级数都收敛. 后来狄利克雷给出了傅里叶级数敛散性判别法. 在狄利克雷的研究工作之后, 人们以为任何连续函数的傅里叶级数都收敛到该函数, 但人们在 1873 年给出了例子, 从这些例子可以看出一个连续函数的傅里叶级数可以在一个无穷点集上发散, 而且后来人们还举出了可积函数的傅里叶级数处处发散的例子. 在傅里叶级数的敛散性这个问题上直到近代还有一些重要的研究工作, 由于这些工作的介绍要用到实变函数课程中的知识, 我们在这里不再赘述。

\section{傅里叶级数的其他收敛性}

在上节中我们主要讨论了函数 $f(x)$ 的傅里叶级数的点收敛性. 在本节中, 我们将讨论傅里叶级数有关的其他一些收敛性. 不失一般性, 在本节中我们假定 $f(x)$ 是周期为 $2\pi$ 的函数. 读者不难将本节的结果推广到 $f(x)$ 是以 $2T$ 为周期的函数的情形。

\subsection{连续函数的三角多项式一致逼近}

我们在幂级数理论中已经证明了 $[a,b]$ 上连续函数必可由多项式一致逼近. 在本小节中, 我们证明以 $2\pi$ 为周期的连续函数必可由三角多项式一致逼近, 即证明

\noindent\textbf{定理 12.3.1 (魏尔斯特拉斯第二逼近定理)}\quad 设 $f(x)$ 是以 $2\pi$ 为周期的连续函数, 则存在
\[
  T_n(x)=\frac{\alpha_0}{2}+\sum_{k=1}^n(\alpha_k\cos kx+\beta_k\sin kx)\quad(n=1,2,\cdots),
\]
使得 $\forall\varepsilon>0$, 存在 $N\in\mathbb N$, 当 $n>N$ 时, 对一切 $x\in(-\infty,+\infty)$, 有
\[
  |T_n(x)-f(x)|<\varepsilon.
\]

为了证明定理 12.3.1, 我们需要做些准备工作. 从上节我们介绍的傅里叶级数敛散性研究的历史中可知, 一般来说, 定理 12.3.1 中的三角多项式序列不能取 $f(x)$ 的傅里叶级数的部分和序列. 因此我们必须另找其他的三角多项式序列. 我们回忆一下, 对一个序列 $\{x_n\}$ 来说, 若 $\lim\limits_{n\to\infty}x_n=a$, 则必有
\[
  \lim_{n\to\infty}\frac{x_1+\cdots+x_n}{n}=a,
\]
而反之一般不真. 因此 $\left\{y_n=\dfrac{x_1+\cdots+x_n}{n}\right\}$ 比 $\{x_n\}$ 来说有可能具有更好的收敛性。

设 $f(x)$ 的傅里叶级数为
\[
  \frac{a_0}{2}+\sum_{n=1}^{+\infty}(a_n\cos nx+b_n\sin nx),
\]
$S_n(x)(n=0,1,2,\cdots)$ 为其部分和序列, 对于 $n=0,1,2,\cdots$, 我们记
\[
  S_n^*(x)=\frac{S_0(x)+S_1(x)+\cdots+S_n(x)}{n+1}.
\]
我们将看到, 对于 $n\in\mathbb N$, $S_n^*(x)$ 也有很好的积分表达式. 事实上, 我们有
\[
\begin{aligned}
  S_n^*(x)
  &=\frac{1}{(n+1)\pi}\int_{-\pi}^{\pi}f(x+t)
    \left[
      \frac{\displaystyle\sum_{k=0}^n\sin\left(k+\frac12\right)t}{2\sin\dfrac t2}
    \right]\mathrm dt\\
  &=\frac{1}{(n+1)\pi}\int_{-\pi}^{\pi}f(x+t)
    \left[
      \frac{\displaystyle\sum_{k=0}^n2\sin\dfrac t2\sin\left(k+\frac12\right)t}{4\sin^2\dfrac t2}
    \right]\mathrm dt\\
  &=\frac{1}{(n+1)\pi}\int_{-\pi}^{\pi}f(x+t)
    \left[
      \frac{\displaystyle\sum_{k=0}^n(\cos kt-\cos(k+1)t)}{4\sin^2\dfrac t2}
    \right]\mathrm dt\\
  &=\frac1\pi\int_{-\pi}^{\pi}f(x+t)
    \frac{\sin^2\dfrac{n+1}{2}t}{2(n+1)\sin^2\dfrac t2}\mathrm dt\\
  &=\frac1\pi\int_{-\pi}^{\pi}f(x+t)\Phi_n(t)\mathrm dt,
\end{aligned}
\]
其中
\[
  \Phi_n(t)=\frac{\sin^2\dfrac{n+1}{2}t}{2(n+1)\sin^2\dfrac t2}
\]
称为费叶 (Fejer) 核. 显然对任意的 $t$, 有 $\Phi_n(t)\geqslant0$. 在上述积分表达式中令 $f(x)\equiv1$, 我们推出
\[
  \frac1\pi\int_{-\pi}^{\pi}\Phi_n(t)\mathrm dt=1.
\]

下面我们来证明定理 12.3.1。

\noindent\textbf{定理 12.3.1 的证明}\quad 我们将证明 $f(x)$ 的傅里叶级数部分和序列 $\{S_n(x)\}$ 的算术平均构成的序列 $\{S_n^*(x)\}$ 在 $[-\pi,\pi]$ 上一致收敛于 $f(x)$. 对于任意的 $x\in[-\pi,\pi]$, 我们有下述不等式:
\begin{equation}
\begin{aligned}
  |f(x)-S_n^*(x)|
  &=\left|\frac1\pi\int_{-\pi}^{\pi}\Phi_n(t)[f(x)-f(x+t)]\mathrm dt\right|\\
  &\leqslant\frac1\pi\int_{-\pi}^{\pi}\Phi_n(t)|f(x)-f(x+t)|\mathrm dt.
\end{aligned}
\tag{12.3.1}
\end{equation}
设 $|f(x)|$ 在 $[-\pi,\pi]$ 上的最大值为 $M$. 注意到 $f(x)$ 以 $2\pi$ 为周期, 由 $f(x)$ 的连续性我们推出它在 $(-\infty,+\infty)$ 上一致连续, 从而对于 $\forall\varepsilon>0$, $\exists\delta>0$, 当 $x_1,x_2$ 满足 $|x_1-x_2|<\delta$ 时, 有
\[
  |f(x_1)-f(x_2)|<\frac{\varepsilon}{3}.
\]
将 (12.3.1) 式右边的积分分解成 $I_1+I_2+I_3$, 其中
\[
  I_1=\frac1\pi\int_{-\pi}^{-\delta}\Phi_n(t)|f(x)-f(x+t)|\mathrm dt,
\]
\[
  I_2=\frac1\pi\int_{-\delta}^{\delta}\Phi_n(t)|f(x)-f(x+t)|\mathrm dt,
\]
\[
  I_3=\frac1\pi\int_{\delta}^{\pi}\Phi_n(t)|f(x)-f(x+t)|\mathrm dt.
\]
现在我们来估计上面的每一个积分。

对于 $I_2$, 我们有
\[
\begin{aligned}
  I_2
  &=\frac1\pi\int_{-\delta}^{\delta}\Phi_n(t)|f(x)-f(x+t)|\mathrm dt\\
  &\leqslant\frac{\varepsilon}{3\pi}\int_{-\delta}^{\delta}\Phi_n(t)\mathrm dt\\
  &\leqslant\frac{\varepsilon}{3\pi}\int_{-\pi}^{\pi}\Phi_n(t)\mathrm dt
  =\frac{\varepsilon}{3}.
\end{aligned}
\]

对于 $I_1$, 我们有
\[
\begin{aligned}
  I_1
  &=\frac1\pi\int_{-\pi}^{-\delta}\Phi_n(t)|f(x)-f(x+t)|\mathrm dt\\
  &\leqslant 2M\frac1\pi\int_{-\pi}^{-\delta}
  \frac{\sin^2\dfrac{n+1}{2}t}{2(n+1)\sin^2\dfrac t2}\mathrm dt\\
  &\leqslant 2M\frac1\pi\int_{-\pi}^{-\delta}
  \frac{1}{2(n+1)\sin^2\dfrac\delta2}\mathrm dt\\
  &<\frac{M\pi}{\sin^2\dfrac\delta2}\cdot\frac1{n+1}.
\end{aligned}
\]
因此存在 $N_1>0$, 当 $n>N_1$ 时, 有
\[
  I_1<\frac{\varepsilon}{3}.
\]
同理可证, 存在 $N_2>0$, 当 $n>N_2$ 时, 有
\[
  I_3<\frac{\varepsilon}{3}.
\]
因此, 当 $n>N=\max\{N_1,N_2\}$ 时, 对于 $\forall x\in[-\pi,\pi]$, 有
\[
  |f(x)-S_n^*(x)|<\varepsilon.
\]
注意到对于 $\forall n\in\mathbb N$, $S_n^*(x)$ 仍为一个至多 $n$ 阶的三角多项式 (且与 $S_n(x)$ 的阶相同), 定理 12.3.1 得证. 证毕。

\noindent\textbf{例 12.3.1}\quad 设函数 $f(x)$ 在 $(-\infty,+\infty)$ 上连续且以 $2\pi$ 为周期, 再设 $f(x)$ 的傅里叶级数处处收敛. 证明 $f(x)$ 的傅里叶级数必处处收敛于 $f(x)$。

\noindent\textbf{证明}\quad 设
\[
  f(x)\sim\frac{a_0}{2}+\sum_{n=1}^{+\infty}(a_n\cos nx+b_n\sin nx).
\]
对于 $\forall x\in(-\infty,+\infty)$ 记
\[
  S_0(x)=\frac{a_0}{2},
\]
\[
  S_n(x)=\frac{a_0}{2}+\sum_{k=1}^n(a_k\cos kx+b_k\sin kx)\quad(n=1,2,\cdots),
\]
\[
  S_n^*(x)=\frac1{n+1}\sum_{k=0}^nS_k(x)\quad(n=1,2,\cdots).
\]
由于 $f(x)$ 的傅里叶级数处处收敛, 因此存在 $(-\infty,+\infty)$ 上定义的函数 $g(x)$, 使得对于 $\forall x\in(-\infty,+\infty)$, 有
\[
  \lim_{n\to\infty}S_n(x)=g(x).
\]
因此, 对于 $\forall x\in(-\infty,+\infty)$, 有
\[
  \lim_{n\to\infty}S_n^*(x)=g(x).
\]
再根据定理 12.3.1, 我们有 $S_n^*(x)\to f(x)(x\in(-\infty,+\infty))$, 从而有 $f(x)\equiv g(x)(x\in(-\infty,+\infty))$, 即
\[
  f(x)=\frac{a_0}{2}+\sum_{n=1}^{+\infty}(a_n\cos nx+b_n\sin nx),\quad x\in(-\infty,+\infty).
\]

\subsection{傅里叶级数的均方收敛}

我们已经对函数序列 (函数项级数) 的点收敛及一致收敛有了许多研究. 对于傅里叶级数来说, 由于它与积分密切相关, 我们可以引入另一种意义下的收敛---均方收敛。

为了讨论均方收敛, 我们必须对所涉及的函数附加一个稍强的条件. 设函数 $f(x)$ 在区间 $[a,b]$ 上除了有限个瑕点外有定义, 在 $[a,b]$ 上的任何不包含瑕点的子区间 $[c,d]$ 上可积, 并且 $f^2(x)$ 在 $[a,b]$ 上的瑕积分是收敛的, 则我们称 $f(x)$ 在 $[a,b]$ 上平方可积. 特别地, 当 $f(x)$ 在 $[a,b]$ 上可积时, $f(x)$ 在 $[a,b]$ 上必平方可积。

\noindent\textbf{定义 12.3.1}\quad 设函数 $f(x),f_n(x)(n=1,2,\cdots)$ 在区间 $[a,b]$ 上平方可积, 并且满足
\[
  \lim_{n\to\infty}\int_a^b[f_n(x)-f(x)]^2\mathrm dx=0,
\]
则称函数序列 $\{f_n(x)\}$ 在 $[a,b]$ 上均方收敛于 $f(x)$。

函数序列的这种均方收敛性在一些数学分支的研究中起着重要的作用, 在傅里叶级数理论中也具有重要的理论和应用价值. 下面我们先来研究一下平方可积函数类的一些初等性质。

由于对区间 $[a,b]$ 上任意的函数 $f(x)$ 总成立
\[
  |f(x)|\leqslant\frac{f^2(x)+1}{2}.
\]
因此, 当 $f(x)$ 在 $[a,b]$ 上平方可积时, $f(x)$ 在 $[a,b]$ 上必定绝对可积. 这说明 $f(x)$ 平方可积的条件要比 $f(x)$ 绝对可积的条件强一些. 对于区间 $[a,b]$ 上平方可积的函数 $f(x),g(x)$, 我们有以下的结论:

(1) $f(x)g(x)$ 在 $[a,b]$ 上绝对可积;

(2) $f(x)+g(x)$ 在 $[a,b]$ 上平方可积。

事实上, 对于 $\forall x\in[a,b]$, 我们有
\[
  |f(x)g(x)|\leqslant\frac{f^2(x)+g^2(x)}{2},
\]
因此 (1) 成立。

记
\[
\begin{aligned}
  \varphi(t)
  &=\int_a^b(t|f|+|g|)^2\mathrm dx\\
  &=t^2\int_a^bf^2(x)\mathrm dx
    +2t\int_a^b|f(x)g(x)|\mathrm dx
    +\int_a^bg^2(x)\mathrm dx.
\end{aligned}
\]
则对于 $\forall t\in\mathbb R$, 有 $\varphi(t)\geqslant0$. 由此我们推知上式右边关于 $t$ 的一元二次方程的判别式 $\Delta\leqslant0$, 即
\[
  \int_a^b|f(x)g(x)|\mathrm dx
  \leqslant
  \left[\int_a^bf^2(x)\mathrm dx\right]^{\frac12}
  \left[\int_a^bg^2(x)\mathrm dx\right]^{\frac12}.
\]
由此不等式我们也可推出 $f(x)g(x)$ 在 $[a,b]$ 上的绝对可积性. 利用这个不等式, 我们还有
\[
\begin{aligned}
  \int_a^b[f(x)+g(x)]^2\mathrm dx
  &=\int_a^bf^2(x)\mathrm dx+2\int_a^bf(x)g(x)\mathrm dx+\int_a^bg^2(x)\mathrm dx\\
  &\leqslant\int_a^bf^2(x)\mathrm dx
    +2\sqrt{\int_a^bf^2(x)\mathrm dx}\sqrt{\int_a^bg^2(x)\mathrm dx}
    +\int_a^bg^2(x)\mathrm dx\\
  &=\left[
    \sqrt{\int_a^bf^2(x)\mathrm dx}
    +\sqrt{\int_a^bg^2(x)\mathrm dx}
  \right]^2.
\end{aligned}
\]
由此得到
\[
  \left\{\int_a^b[f(x)+g(x)]^2\mathrm dx\right\}^{\frac12}
  \leqslant
  \left[\int_a^bf^2(x)\mathrm dx\right]^{\frac12}
  +
  \left[\int_a^bg^2(x)\mathrm dx\right]^{\frac12}.
\]
上面的不等式也称为柯西不等式, 从它可以直接看出 $f(x)+g(x)$ 是平方可积的. (2) 得证。

现在我们来考虑函数 $f(x)$ 的傅里叶级数的均方收敛性问题. 显然, 当 $f(x)$ 在 $[-\pi,\pi]$ 上可积或有瑕点时平方可积时, 对任意的三角多项式
\[
  T_n(x)=\frac{\alpha_0}{2}+\sum_{k=1}^n(\alpha_k\cos kx+\beta_k\sin kx),
\]
$f(x)-T_n(x)$ 也在 $[-\pi,\pi]$ 上可积或有瑕点时平方可积. 设 $f(x)$ 的傅里叶级数为
\[
  \frac{a_0}{2}+\sum_{n=1}^{+\infty}(a_n\cos nx+b_n\sin nx),
\]
令
\[
  \Delta_n=\left\{\frac1{2\pi}\int_{-\pi}^{\pi}[f(x)-T_n(x)]^2\mathrm dx\right\}^{\frac12},
\]
则我们称 $\Delta_n$ 为 $f(x)-T_n(x)$ 的平均偏差. 下面我们来计算 $\Delta_n^2$:
\[
\begin{aligned}
  \Delta_n^2
  &=\frac1{2\pi}\int_{-\pi}^{\pi}[f(x)-T_n(x)]^2\mathrm dx\\
  &=\frac1{2\pi}\left[
    \int_{-\pi}^{\pi}f^2(x)\mathrm dx
    -2\int_{-\pi}^{\pi}f(x)T_n(x)\mathrm dx
    +\int_{-\pi}^{\pi}T_n^2(x)\mathrm dx
  \right]\\
  &=\frac1{2\pi}\int_{-\pi}^{\pi}f^2(x)\mathrm dx
    -\frac1\pi\int_{-\pi}^{\pi}f(x)
    \left[\frac{\alpha_0}{2}+\sum_{k=1}^n(\alpha_k\cos kx+\beta_k\sin kx)\right]\mathrm dx\\
  &\quad+\frac1{2\pi}\int_{-\pi}^{\pi}
    \left[\frac{\alpha_0}{2}+\sum_{k=1}^n(\alpha_k\cos kx+\beta_k\sin kx)\right]^2\mathrm dx\\
  &=\frac1{2\pi}\int_{-\pi}^{\pi}f^2(x)\mathrm dx
    -\frac{\alpha_0a_0}{2}
    -\sum_{k=1}^n(\alpha_ka_k+\beta_kb_k)
    +\frac{\alpha_0^2}{4}
    +\frac12\sum_{k=1}^n(\alpha_k^2+\beta_k^2)\\
  &=\frac1{2\pi}\int_{-\pi}^{\pi}f^2(x)\mathrm dx
    +\frac12\left[
      -\alpha_0a_0+\frac{\alpha_0^2}{2}
      -2\sum_{k=1}^n(\alpha_ka_k+\beta_kb_k)
      +\sum_{k=1}^n(\alpha_k^2+\beta_k^2)
    \right]\\
  &=\frac1{2\pi}\int_{-\pi}^{\pi}f^2(x)\mathrm dx
    +\frac12\left[
      \frac{(\alpha_0-a_0)^2}{2}
      +\sum_{k=1}^n\{(\alpha_k-a_k)^2+(\beta_k-b_k)^2\}
    \right]\\
  &\quad-\frac{a_0^2}{4}
    -\frac12\sum_{k=1}^n(a_k^2+b_k^2).
\end{aligned}
\]
从上面的等式我们可以得出以下一系列的结论。

\noindent\textbf{定理 12.3.2 (傅里叶级数最佳逼近)}\quad 设函数 $f(x)$ 在 $[-\pi,\pi]$ 上可积或有瑕点时平方可积, 并设其傅里叶级数的部分和序列为 $\{S_n(x)\}$, 则对任何 $n\in\mathbb N$ 阶三角多项式 $T_n(x)$, 成立
\[
  \frac1{2\pi}\int_{-\pi}^{\pi}[f(x)-S_n(x)]^2\mathrm dx
  \leqslant
  \frac1{2\pi}\int_{-\pi}^{\pi}[f(x)-T_n(x)]^2\mathrm dx,
\]
并且当且仅当 $S_n(x)\equiv T_n(x)$ 时上面不等式中等号成立。

\noindent\textbf{注}\quad 从上述分析我们还知道, 在定理 12.3.2 的条件下, 对于 $f(x)$ 的傅里叶级数部分和序列有以下性质: 当 $m<n$ 时, 有
\[
  \int_{-\pi}^{\pi}[f(x)-S_n(x)]^2\mathrm dx
  \leqslant
  \int_{-\pi}^{\pi}[f(x)-S_m(x)]^2\mathrm dx.
\]

由定理 12.3.2 可以解决我们在本小节的中心问题:

\noindent\textbf{定理 12.3.3}\quad 设函数 $f(x)$ 在 $[-\pi,\pi]$ 上可积或有瑕点时平方可积, 则对 $f(x)$ 的傅里叶级数部分和序列 $\{S_n(x)\}$, 有
\[
  \lim_{n\to\infty}\int_{-\pi}^{\pi}[f(x)-S_n(x)]^2\mathrm dx=0.
\]

\noindent\textbf{证明}\quad 先设 $f(x)$ 在 $[-\pi,\pi]$ 上连续且以 $2\pi$ 为周期, 由魏尔斯特拉斯第二逼近定理, 对于 $\forall\varepsilon>0$, 存在三角多项式 $T_{n_0}(x)$, 使得对 $\forall x\in[-\pi,\pi]$, 有
\[
  |f(x)-T_{n_0}(x)|<\sqrt{\frac{\varepsilon}{2\pi}}.
\]
因此我们有
\[
  \int_{-\pi}^{\pi}[f(x)-T_{n_0}(x)]^2\mathrm dx<\varepsilon.
\]
现设 $f(x)$ 在 $[-\pi,\pi]$ 上可积. 我们知道存在连续函数 (可以是分段线性函数 (见例 7.3.1 的注)) $g(x)$, 使得
\[
  \int_{-\pi}^{\pi}[f(x)-g(x)]^2\mathrm dx<\frac{\varepsilon}{4}.
\]
在 $x=\pm\pi$ 的充分小的邻域内对 $g(x)$ 的值进行修改, 我们可以假定满足上述不等式的 $g(x)$ 同时满足 $g(-\pi)=g(\pi)$。

由上面所证, 存在 $n_0\in\mathbb N$, 使得
\[
  \int_{-\pi}^{\pi}[g(x)-T_{n_0}(x)]^2\mathrm dx<\frac{\varepsilon}{4}.
\]
因此
\[
\begin{aligned}
  \int_{-\pi}^{\pi}[f(x)-T_{n_0}(x)]^2\mathrm dx
  &\leqslant
  2\left\{
    \int_{-\pi}^{\pi}[f(x)-g(x)]^2\mathrm dx
    +\int_{-\pi}^{\pi}[g(x)-T_{n_0}(x)]^2\mathrm dx
  \right\}\\
  &<\varepsilon.
\end{aligned}
\]
最后设 $f(x)$ 有瑕点且在 $[-\pi,\pi]$ 上平方可积. 此时不妨设 $f(x)$ 仅以 $x=\pi$ 为瑕点. 由假设, $\exists\delta>0$, 使得
\[
  \int_{\pi-\delta}^{\pi}f^2(x)\mathrm dx<\frac{\varepsilon}{8}
  \quad\text{及}\quad
  \int_{-\pi}^{-\pi+\delta}f^2(x)\mathrm dx<\frac{\varepsilon}{8}.
\]
令
\[
  \widetilde f(x)=
  \begin{cases}
    0, & x\in(\pi-\delta,\pi],\\
    f(x), & x\in[-\pi+\delta,\pi-\delta],\\
    0, & x\in[-\pi,-\pi+\delta).
  \end{cases}
\]
显然 $\widetilde f(x)$ 在 $[-\pi,\pi]$ 上可积, 且 $\widetilde f(-\pi)=\widetilde f(\pi)$, 从而存在三角多项式 $T_{n_0}(x)$, 使得
\[
  \int_{-\pi}^{\pi}[\widetilde f(x)-T_{n_0}(x)]^2\mathrm dx<\frac{\varepsilon}{4}.
\]
因此
\[
\begin{aligned}
  \int_{-\pi}^{\pi}[f(x)-T_{n_0}(x)]^2\mathrm dx
  &\leqslant
  2\int_{-\pi}^{\pi}[\widetilde f(x)-f(x)]^2\mathrm dx
  +2\int_{-\pi}^{\pi}[\widetilde f(x)-T_{n_0}(x)]^2\mathrm dx\\
  &\leqslant
  2\int_{\pi-\delta}^{\pi}f^2(x)\mathrm dx
  +2\int_{-\pi}^{-\pi+\delta}f^2(x)\mathrm dx
  +\frac{\varepsilon}{2}\\
  &<\frac{\varepsilon}{4}+\frac{\varepsilon}{4}+\frac{\varepsilon}{2}
  =\varepsilon.
\end{aligned}
\]
综合上述的证明我们知, 若 $f(x)$ 在 $[-\pi,\pi]$ 上可积或有瑕点时平方可积, 则存在三角多项式 $T_{n_0}(x)$, 使得
\[
  \int_{-\pi}^{\pi}[f(x)-T_{n_0}(x)]^2\mathrm dx<\varepsilon.
\]
由定理 12.3.2, 对于 $f(x)$ 的傅里叶级数部分和 $S_{n_0}(x)$, 有
\[
  \int_{-\pi}^{\pi}[f(x)-S_{n_0}(x)]^2\mathrm dx
  \leqslant
  \int_{-\pi}^{\pi}[f(x)-T_{n_0}(x)]^2\mathrm dx<\varepsilon.
\]
因此当 $n>N(=n_0)$ 时, 有
\[
  \int_{-\pi}^{\pi}[f(x)-S_n(x)]^2\mathrm dx
  \leqslant
  \int_{-\pi}^{\pi}[f(x)-S_{n_0}(x)]^2\mathrm dx<\varepsilon.
\]
证毕。

由定理 12.3.3 我们可以推出下面的帕塞瓦尔 (Parseval) 等式。

\noindent\textbf{定理 12.3.4 (帕塞瓦尔等式)}\quad 设函数 $f(x)$ 在 $[-\pi,\pi]$ 上可积或有瑕点时平方可积, 则有
\[
  \frac{a_0^2}{2}+\sum_{n=1}^{+\infty}(a_n^2+b_n^2)
  =
  \frac1\pi\int_{-\pi}^{\pi}f^2(x)\mathrm dx.
\]

\noindent\textbf{证明}\quad 由定理 12.3.3, 对于 $\forall\varepsilon>0$, $\exists N\in\mathbb N$, 当 $n>N$ 时, 有
\[
  \frac1\pi\int_{-\pi}^{\pi}[f(x)-S_n(x)]^2\mathrm dx<\varepsilon,
\]
其中 $\{S_n(x)\}$ 是 $f(x)$ 的傅里叶级数部分和序列. 因此我们有
\[
\begin{aligned}
  0
  &\leqslant\frac1\pi\int_{-\pi}^{\pi}f^2(x)\mathrm dx
  -\left[\frac{a_0^2}{2}+\sum_{k=1}^n(a_k^2+b_k^2)\right]\\
  &=\frac1\pi\int_{-\pi}^{\pi}[f(x)-S_n(x)]^2\mathrm dx
  <\varepsilon.
\end{aligned}
\]
证毕。

从帕塞瓦尔等式我们可以知道, 并不是任意一个三角级数都是某个平方可积函数的傅里叶级数. 例如, 对于三角级数
\[
  \sum_{n=1}^{+\infty}\frac{\cos nx}{\sqrt n},
\]
由于
\[
  \sum_{n=1}^{+\infty}\frac1{(\sqrt n)^2}=+\infty,
\]
由帕塞瓦尔等式, 它不可能是某个平方可积函数 $f(x)$ 的傅里叶级数。

\noindent\textbf{注}\quad 设 $f(x)$ 与 $g(x)$ 均是以 $2\pi$ 为周期的函数且在 $[-\pi,\pi]$ 上平方可积, 并设 $f(x)$ 与 $g(x)$ 的傅里叶级数分别为
\[
  f(x)\sim\frac{a_0}{2}+\sum_{n=1}^{+\infty}(a_n\cos nx+b_n\sin nx),
\]
\[
  g(x)\sim\frac{\alpha_0}{2}+\sum_{n=1}^{+\infty}(\alpha_n\cos nx+\beta_n\sin nx),
\]
则成立
\begin{equation}
  \frac1\pi\int_{-\pi}^{\pi}f(x)g(x)\mathrm dx
  =
  \frac{a_0\alpha_0}{2}+\sum_{n=1}^{+\infty}(a_n\alpha_n+b_n\beta_n),
  \tag{12.3.2}
\end{equation}
事实上, 从 $f+g$ 与 $f-g$ 的帕塞瓦尔等式相减即可得到 (12.3.2). (12.3.2) 也称为广义帕塞瓦尔等式。

\subsection{傅里叶级数的一致收敛性}

利用帕塞瓦尔等式, 我们可以容易解决傅里叶级数的一致收敛、逐项求导以及逐项积分的问题。

\noindent\textbf{定理 12.3.5 (傅里叶级数的一致收敛性)}\quad 设函数 $f(x)$ 以 $2\pi$ 为周期, 在 $[-\pi,\pi]$ 上可导, 且 $f'(x)$ 在 $[-\pi,\pi]$ 上可积, 则 $f(x)$ 的傅里叶级数在 $(-\infty,+\infty)$ 上一致收敛到 $f(x)$。

\noindent\textbf{证明}\quad 设 $f(x)$ 的傅里叶级数为
\[
  f(x)\sim\frac{a_0}{2}+\sum_{n=1}^{+\infty}(a_n\cos nx+b_n\sin nx),
\]
$f'(x)$ 的傅里叶级数为
\[
  f'(x)\sim\frac{a_0'}{2}+\sum_{n=1}^{+\infty}(a_n'\cos nx+b_n'\sin nx).
\]
我们有
\[
  a_0'=\frac1\pi\int_{-\pi}^{\pi}f'(x)\mathrm dx
  =\frac1\pi f(x)\bigg|_{-\pi}^{\pi}=0,
\]
\[
\begin{aligned}
  a_n'
  &=\frac1\pi\int_{-\pi}^{\pi}f'(x)\cos nx\mathrm dx\\
  &=\frac1\pi f(x)\cos nx\bigg|_{-\pi}^{\pi}
    +\frac n\pi\int_{-\pi}^{\pi}f(x)\sin nx\mathrm dx\\
  &=nb_n,\quad n=1,2,\cdots.
\end{aligned}
\]
同理我们有
\[
  b_n'=-na_n,\quad n=1,2,\cdots.
\]
因此对于 $\forall N\in\mathbb N$, 有
\[
\begin{aligned}
  \sum_{n=1}^N(|a_n|+|b_n|)
  &=\sum_{n=1}^N\frac{|a_n'|+|b_n'|}{n}\\
  &\leqslant
  \left[\sum_{n=1}^N(a_n'^2+b_n'^2)\right]^{\frac12}
  \left(\sum_{n=1}^N\frac2{n^2}\right)^{\frac12}\\
  &<\left(\frac{\pi^2}{3}\right)^{\frac12}
  \left[\frac1\pi\int_{-\pi}^{\pi}(f'(x))^2\mathrm dx\right]^{\frac12}.
\end{aligned}
\]
因此 $f(x)$ 的傅里叶级数在 $(-\infty,+\infty)$ 上绝对一致收敛. 由于 $f(x)$ 处处可导, $f(x)$ 的傅里叶级数处处收敛到 $f(x)$, 因此 $f(x)$ 的傅里叶级数在 $(-\infty,+\infty)$ 上一致收敛到 $f(x)$. 证毕。

\noindent\textbf{定理 12.3.6 (傅里叶级数逐项微分定理)}\quad 设函数 $f(x)$ 以 $2\pi$ 为周期, 对于 $x\in[-\pi,\pi]$, $f''(x)$ 存在且在 $[-\pi,\pi]$ 上可积, 再设 $f(x)$ 的傅里叶级数为
\[
  f(x)=\frac{a_0}{2}+\sum_{n=1}^{+\infty}(a_n\cos nx+b_n\sin nx),\quad x\in(-\infty,+\infty),
\]
则
\[
  f'(x)=\sum_{n=1}^{+\infty}(nb_n\cos nx-na_n\sin nx),\quad x\in(-\infty,+\infty).
\]

\noindent\textbf{证明}\quad 直接应用定理 12.3.5 及函数项级数逐项微分定理即得. 证毕。

对于 $f(x)$ 的傅里叶级数的逐项积分问题, 我们在相当弱的条件下即能得到相应的结果。

\noindent\textbf{定理 12.3.7 (傅里叶级数逐项积分定理)}\quad 设函数 $f(x)$ 在 $[0,2\pi]$ 上可积, 且以 $2\pi$ 为周期, 再设
\[
  f(x)\sim\frac{a_0}{2}+\sum_{n=1}^{+\infty}(a_n\cos nx+b_n\sin nx),
\]
则
\[
  \int_0^xf(t)\mathrm dt
  =
  \frac{a_0}{2}x+\sum_{n=1}^{+\infty}
  \left[
    \frac{a_n}{n}\sin nx+\frac{b_n(1-\cos nx)}{n}
  \right]
\]
在 $[0,2\pi]$ 成立。

\noindent\textbf{证明}\quad 对于 $x\in[0,2\pi]$, 作
\[
  g(t)=
  \begin{cases}
    \dfrac{\pi}{2}, & t=0,t=x,\\
    \pi, & t\in(0,x),\\
    0, & t\in(x,2\pi).
  \end{cases}
\]
设 $g(t)$ 的傅里叶级数为
\[
  \frac{\alpha_0}{2}+\sum_{n=1}^{+\infty}(\alpha_n\cos nt+\beta_n\sin nt),
\]
则由计算直接得出
\[
  \alpha_0=x,\quad
  \alpha_n=\frac{\sin nx}{n}\quad(n=1,2,\cdots),\quad
  \beta_n=\frac{1-\cos nx}{n}\quad(n=1,2,\cdots).
\]
由关于 $f(t)$ 与 $g(t)$ 的帕塞瓦尔等式即得
\[
  \frac1\pi\int_0^x\pi f(t)\mathrm dt
  =
  \frac{a_0}{2}x
  +\sum_{n=1}^{+\infty}
  \left[
    \frac{a_n}{n}\sin nx+\frac{b_n(1-\cos nx)}{n}
  \right],
  \quad x\in[0,2\pi].
\]
证毕。

\noindent\textbf{例 12.3.2}\quad 求函数 $f(x)=\dfrac{x^2}{2}-\pi x(x\in[0,2\pi])$ 的傅里叶级数, 并求其和函数。

\noindent\textbf{解}\quad 由例 12.2.2 有
\[
  \frac{\pi-x}{2}\sim\sum_{n=1}^{+\infty}\frac{\sin nx}{n},
\]
因此由定理 12.3.7, 对于 $\forall x\in[0,2\pi]$, 有
\[
  \int_0^x\frac{\pi-t}{2}\mathrm dt
  =
  \sum_{n=1}^{+\infty}\int_0^x\frac{\sin nt}{n}\mathrm dt,
\]
即
\[
  \frac{x^2}{4}-\frac{\pi}{2}x
  =
  -\sum_{n=1}^{+\infty}\frac1{n^2}
  +\sum_{n=1}^{+\infty}\frac{\cos nx}{n^2},\quad x\in[0,2\pi].
\]
所以
\[
  \frac{x^2}{2}-\pi x
  =
  -\frac{\pi^2}{3}+\sum_{n=1}^{+\infty}\frac{2\cos nx}{n^2},\quad x\in[0,2\pi].
\]

\section*{习题十二}

1. 求下列周期为 $2\pi$ 的函数的傅里叶级数:

(1) $f(x)=|x|,\ -\pi\leqslant x<\pi$;

(2) $f(x)=e^{\alpha x},\ -\pi\leqslant x<\pi$;

(3) $f(x)=x\cos x,\ -\pi\leqslant x<\pi$;

(4) $f(x)=\cos^2x,\ -\pi\leqslant x<\pi$;

(5)
\[
  f(x)=
  \begin{cases}
    \dfrac{\pi}{4}, & 0<x<\pi,\\
    -\dfrac{\pi}{4}, & -\pi\leqslant x\leqslant0;
  \end{cases}
\]

(6) $f(x)=\sin^4x$;

(7)
\[
  f(x)=\frac{1-r^2}{1-2r\cos x+r^2}\quad(|r|<1);
\]

(8) $f(x)=\ln\left|\sin\dfrac x2\right|$。

\end{document}
